\chapter{Symmetrie, Entartung und Drehimpuls-Leiteroperatoren}
\providecommand{\ket}[1]{\left\lvert#1\right\rangle}

\section{Symmetrien und entartete Energieniveaus}

Eine Symmetrie ist eine Operation, die die physikalische Beschreibung eines Systems nicht verändert. Sie verändert weder den Hamiltonoperator noch die Schr\"odinger-Gleichung oder die möglichen Energien. Wird die Operation durch einen unitären Operator \(S\) dargestellt, bedeutet Invarianz
\begin{equation}
S H S^\dagger=H,
\end{equation}
oder äquivalent dazu
\begin{equation}
[S,H]=0.
\end{equation}

Eine Translation ist ein Beispiel. Verschiebt man ein Atom von einem Ort zu einem anderen, ändert sich seine Energie nicht. Das gilt für jeden Zustand des Atoms, nicht nur für den Grundzustand. Ein weiteres Beispiel ist die Rotation, um die es hier hauptsächlich gehen wird.

Auch die Vertauschung zweier identischer Teilchen ist eine Symmetrie. Zwei Elektronen Elektron eins und Elektron zwei oder Harry und Fred zu nennen, macht sie nicht verschieden. Eine Vertauschung der Bezeichnungen kann die Eigenschaften des Atoms nicht ändern.

Ein Kristall besitzt eine weitere Art von Translationssymmetrie. Den ganzen Kristall durch den Raum zu verschieben, ist eine gewöhnliche Translation. Bei einem unendlichen Kristall ergibt eine Translation um einen Gitterabstand jedoch wieder dieselbe physikalische Anordnung. Äquivalent dazu lassen sich die Gitterplätze umbenennen, indem man jeden Platzindex um eins verschiebt. Bei einem endlichen Kristall verhindern die Ränder diese exakte Aussage; weit entfernt von den Enden eines sehr langen Kristalls ist sie jedoch eine gute Symmetrie.

Die Homogenität des Raums---die Tatsache, dass jeder Ort jedem anderen gleicht---ist Translationssymmetrie. Die Isotropie des Raums---die Tatsache, dass jede Richtung jeder anderen gleicht---ist Rotationssymmetrie.

Entartung hängt mit Symmetrie zusammen, ist aber nicht dasselbe. Ein Energieniveau ist entartet, wenn mehr als ein Zustand genau dieselbe Energie hat. Bei einem generischen Hamiltonoperator mit diskretem Spektrum wäre die exakte Gleichheit zweier nicht miteinander verbundener Niveaus ein bemerkenswerter Zufall. Übereinstimmung auf einen Teil in \(10^{15}\) ist noch keine exakte Gleichheit. Exakte Entartungen weisen gewöhnlich auf eine Symmetrie hin.

Symmetrie kann manchmal Entartung zur Folge haben, aber nicht immer. Das erste Beispiel zeigt sowohl, warum das so sein kann, als auch, weshalb eine einzelne Symmetrie manchmal nicht ausreicht.

\section{Ein Teilchen auf einem Kreis}

Betrachten wir eine Perle, die sich auf einem kreisförmigen Draht bewegen kann. Ihre Position ließe sich durch die kartesischen Koordinaten \(x\) und \(y\) beschreiben, doch eine Winkelkoordinate genügt:
\begin{equation}
\psi=\psi(\theta).
\end{equation}
Die Wahrscheinlichkeitsdichte der Position auf dem Kreis ist \(\psi^*(\theta)\psi(\theta)\).

Eine aktive Drehung gegen den Uhrzeigersinn um einen kleinen Winkel \(\epsilon\) transformiert
\begin{equation}
\psi(\theta)\longrightarrow\psi(\theta-\epsilon).
\end{equation}
Daher gilt
\begin{align}
\delta\psi(\theta)
&=\psi(\theta-\epsilon)-\psi(\theta)\\
&=-\epsilon\frac{d\psi}{d\theta}
+\mathcal O(\epsilon^2).
\end{align}

Der Operator
\begin{equation}
L_z=-i\hbar\frac{d}{d\theta}
\end{equation}
ist hermitesch auf dem Raum der periodischen Wellenfunktionen. Er ist das Winkelanalogon von
\begin{equation}
p_x=-i\hbar\frac{d}{dx}.
\end{equation}
Der Faktor \(i\) ist für die Hermitizität entscheidend. In Bezug auf \(L_z\) gilt
\begin{equation}
\delta\psi
=-\frac{i\epsilon}{\hbar}L_z\psi
+\mathcal O(\epsilon^2),
\end{equation}
Der infinitesimale Rotationsoperator lautet somit
\begin{equation}
R(\epsilon)
=1-\frac{i\epsilon}{\hbar}L_z
+\mathcal O(\epsilon^2).
\end{equation}
Damit ist \(L_z\) als Erzeuger von Drehungen um die zur Kreisebene senkrechte Achse identifiziert.

\begin{classroomqa}
\lecturetimestamp{00:11:19}
\audiencequestion{Müsste nicht einer der Faktoren mit \(i\) das entgegengesetzte Vorzeichen haben?}
\lecturerresponse{Eine Überprüfung der Vorzeichen ergibt die obige Beziehung. Mit der Konvention \(\psi(\theta)\to\psi(\theta-\epsilon)\) ist die Änderung \(-\epsilon\,d\psi/d\theta\), was gleich \(-i\epsilon L_z\psi/\hbar\) ist.}
\end{classroomqa}

Bestimmen wir nun die Eigenzustände des Drehimpulses:
\begin{equation}
L_z\psi=M\psi.
\end{equation}
Einsetzen des Differentialoperators ergibt
\begin{equation}
-i\hbar\frac{d\psi}{d\theta}=M\psi,
\end{equation}
Deren Lösungen haben die Form
\begin{equation}
\psi(\theta)=C\exp\left(\frac{iM\theta}{\hbar}\right).
\end{equation}

Es gibt eine zusätzliche Bedingung, die bei der Bewegung auf einer unendlichen Geraden nicht auftritt. Die Punkte \(\theta=0\) und \(\theta=2\pi\) sind derselbe Punkt. Daher muss eine eindeutige orbitale Wellenfunktion erfüllen:
\begin{equation}
\psi(\theta+2\pi)=\psi(\theta).
\end{equation}
Für die obige Eigenfunktion erfordert das
\begin{equation}
\exp\left(\frac{2\pi iM}{\hbar}\right)=1.
\end{equation}
Folglich gilt
\begin{equation}
M=m\hbar,
\qquad
m\in\mathbb Z.
\end{equation}
Die möglichen orbitalen Drehimpulse auf dem Kreis sind daher
\begin{equation}
\ldots,-2\hbar,-\hbar,0,\hbar,2\hbar,\ldots.
\end{equation}

Muss ein Zustand mit Drehimpuls \(m\hbar\) dieselbe Energie haben wie ein Zustand mit Drehimpuls \(-m\hbar\)? Positive und negative Werte von \(m\) beschreiben entgegengesetzte Umlaufrichtungen. Die Frage ist, ob die Rotationssymmetrie allein verlangt, dass diese beiden Bewegungen dieselbe Energie haben.

\begin{classroomqa}
\lecturetimestamp{00:22:51}
\audiencequestion{Bedeutet die Verwendung von \(\psi(\theta-\epsilon)\), dass sich das Teilchen im Uhrzeigersinn bewegt?}
\lecturerresponse{Nein. Diese Konvention legt lediglich fest, welche Richtung als positive Drehung bezeichnet wird. Sie definiert die Vorzeichen von Winkelposition und Drehimpuls, bestimmt aber nicht die tatsächliche Bewegungsrichtung des Teilchens.}
\end{classroomqa}

Ein Magnetfeld senkrecht zur Kreisebene liefert ein Gegenbeispiel. Es kann die Energien von positivem und negativem Drehimpuls in entgegengesetzte Richtungen verschieben und dadurch das Paar aufspalten. Das Vorzeichen der Verschiebung hängt von den Konventionen für Ladung und Orientierung ab; entscheidend ist jedoch die Aufspaltung.

Ein solches Feld erhält weiterhin die Drehungen um die senkrechte Achse. Es zeigt daher, dass die Rotationssymmetrie allein nicht verlangt, dass
\begin{equation}
E(m)=E(-m).
\end{equation}
Ohne die Aufspaltung können \(m=1\) und \(m=-1\) ein entartetes Paar bilden, ebenso \(m=2\) und \(m=-2\). Der Zustand \(m=0\) hat keinen verschiedenen Partner dieser Art.

\section{Spiegelungssymmetrie und das Paar \texorpdfstring{\(m,-m\)}{m, -m}}

Fügen wir eine diskrete Spiegelungssymmetrie hinzu. Sei \(\mathcal M\) die Spiegelung an einem Durchmesser des Kreises:
\begin{equation}
(\mathcal M\psi)(\theta)=\psi(-\theta).
\end{equation}
Sie kehrt die Umlaufrichtung um:
\begin{equation}
\mathcal M\psi_m=\psi_{-m}.
\end{equation}

Ist die Spiegelung eine Symmetrie, so gilt
\begin{equation}
[\mathcal M,H]=0.
\end{equation}
Nehmen wir an,
\begin{equation}
H\ket{m}=E_m\ket{m}.
\end{equation}
Die Anwendung von \(\mathcal M\) ergibt
\begin{align}
H\mathcal M\ket{m}
&=\mathcal M H\ket{m}\\
&=E_m\mathcal M\ket{m}.
\end{align}
Da \(\mathcal M\ket{m}\) der Zustand mit Drehimpuls \(-m\hbar\) ist, gilt
\begin{equation}
E_{-m}=E_m.
\end{equation}
Die Rotationssymmetrie erhält den Drehimpuls; die Spiegelungssymmetrie verknüpft die beiden entgegengesetzten Werte. Zusammen erzwingen sie die Entartung des Paars.

\begin{classroomqa}
\lecturetimestamp{00:29:54}
\audiencequestion{Da die Moden \(m\) und \(-m\) dieselbe Wahrscheinlichkeitsdichte haben, zeigt das nicht bereits, dass ihre Energien gleich sind?}
\lecturerresponse{Nein. Die Wellenfunktionen \(e^{im\theta}\) und \(e^{-im\theta}\) sind verschieden, obwohl ihre Betragsquadrate übereinstimmen. Die Drehimpuls-Eigenwertgleichung betrifft die Wellenfunktion, nicht nur \(\psi^*\psi\). Die Quantenmechanik erfordert Amplituden, nicht nur Wahrscheinlichkeiten.}
\end{classroomqa}

Das Magnetfeld erhält die Drehungen um seine Achse, bricht aber die oben betrachtete Spiegelungssymmetrie.

\begin{classroomqa}
\lecturetimestamp{00:33:05}
\audiencequestion{Kehrt eine Spiegelungssymmetrie das Magnetfeld um?}
\lecturerresponse{Ja. Ein Magnetfeld ist ein Axialvektor. Bei einer Spiegelung wird ein in die Tafelebene gerichtetes Feld zu einem aus der Tafelebene heraus gerichteten Feld. Äquivalent dazu fließt der Strom in einer Leiterschleife und ihrem Spiegelbild in entgegengesetzten Richtungen, sodass sie entgegengesetzt gerichtete Felder erzeugen.}
\end{classroomqa}

Die Nichtvertauschbarkeit von Drehung und Spiegelung lässt sich direkt erkennen. Auf eine Drehimpuls-Eigenfunktion angewandt gilt
\begin{align}
\mathcal M L_z\psi_m
&=m\hbar\,\psi_{-m},\\
L_z\mathcal M\psi_m
&=-m\hbar\,\psi_{-m}.
\end{align}
Somit gilt
\begin{equation}
[\mathcal M,L_z]\psi_m
=2m\hbar\,\psi_{-m},
\end{equation}
was für \(m\neq0\) nicht null ist. Erst zu drehen und dann zu spiegeln ist nicht dieselbe Operation wie erst zu spiegeln und dann zu drehen.

Wenn beide Transformationen Symmetrien sind, vertauschen beide Operatoren mit \(H\), aber nicht miteinander. Dieses Muster ist für die Entartung verantwortlich.

\begin{classroomqa}
\lecturetimestamp{00:40:28}
\audiencequestion{Erzeugt jedes nichtvertauschende Heisenberg-Paar, etwa \(x\) und \(p\), Entartungen?}
\lecturerresponse{Im Allgemeinen nicht. Die Operatoren müssen außerdem Symmetrien sein, also mit dem Hamiltonoperator vertauschen. Ort und Impuls sind normalerweise nicht beide Erhaltungsgrößen. Würden beide mit \(H\) vertauschen, so würde ihre Nichtvertauschbarkeit Entartung zur Folge haben.}
\end{classroomqa}

\section{Kommutatoralgebren von Symmetrien}

Seien \(A\) und \(B\) hermitesche Erzeuger zweier kontinuierlicher Symmetrien:
\begin{equation}
[A,H]=0,
\qquad
[B,H]=0.
\end{equation}
Nichts in diesen Gleichungen verlangt, dass \(A\) und \(B\) miteinander vertauschen.

Der Kommutator \([A,B]\) ist antihermitesch. Definieren wir daher den hermiteschen Operator
\begin{equation}
C=i[A,B].
\end{equation}
Zahlenfaktoren und ein globales Vorzeichen sind hier unerheblich. Entscheidend ist, dass auch \(C\) mit dem Hamiltonoperator vertauscht. In der Tat gilt
\begin{equation}
[AB,H]
=A[B,H]+[A,H]B=0,
\end{equation}
und entsprechend
\begin{equation}
[BA,H]=0.
\end{equation}
Daher gilt
\begin{equation}
[C,H]
=i\bigl([AB,H]-[BA,H]\bigr)=0.
\end{equation}

Der neue Erzeuger \(C\) kann unabhängig von \(A\) und \(B\) sein oder lediglich eine Linearkombination bereits bekannter Erzeuger darstellen. Ist er neu, fügen wir ihn der Liste hinzu und bilden weitere Kommutatoren. Dieser Vorgang kann sich unbegrenzt fortsetzen, doch oft schließt er sich nach endlich vielen Erzeugern: Jeder weitere Kommutator ist dann eine Linearkombination der bereits vorhandenen. Dieses abgeschlossene System heißt Kommutatoralgebra oder Lie-Algebra.

Die endlichen Symmetrietransformationen bilden eine Gruppe. Ihre infinitesimalen Erzeuger bilden die zugehörige Lie-Algebra.

Entartungen sind wichtig, weil sie Spektren strukturieren. Haben zwei Zustände dieselbe Energie, kann der eine nicht in den anderen übergehen, indem er ein Photon abstrahlt, das ihre Energiedifferenz trägt. In der Atomspektroskopie bestimmen Anordnung und Entartung der Energieniveaus die möglichen Spektrallinien. Dieselbe Überlegung gilt für angeregte Hadronen und ihre Übergänge.

\begin{classroomqa}
\lecturetimestamp{00:49:07}
\audiencequestion{Sind diese Kommutatoren selbst Elemente der Symmetriegruppe?}
\lecturerresponse{Bei den betrachteten Operatoren handelt es sich um Erzeuger. Die Erzeuger und ihre Kommutatoren bilden eine Kommutatoralgebra oder Lie-Algebra. Endliche Symmetrieoperationen, die aus den Erzeugern aufgebaut sind, sind die Gruppenelemente.}
\end{classroomqa}

Wenn alle Symmetrieerzeuger miteinander vertauschen, heißt die Symmetrie abelsch. Sind einige ihrer Kommutatoren ungleich null, ist sie nichtabelsch.

\begin{classroomqa}
\lecturetimestamp{00:50:25}
\audiencequestion{Was ist die grundlegende Definition eines Erzeugers?}
\lecturerresponse{Man nehme eine unitäre Symmetrietransformation, die sehr nahe an der Identität liegt, und schreibe sie als \(1-i\epsilon G/\hbar+\cdots\). Der hermitesche Operator \(G\) ist ihr Erzeuger. Durch wiederholte Anwendung der entsprechenden infinitesimalen Transformation lässt sich eine endliche Transformation aufbauen.}
\end{classroomqa}

\begin{classroomqa}
\lecturetimestamp{00:52:08}
\audiencequestion{Was versteht man unter der Algebra der Symmetrieerzeuger?}
\lecturerresponse{Sie ist eine Menge von Erzeugern, die addiert, subtrahiert, mit gewöhnlichen Zahlen multipliziert und miteinander kommutiert werden können. In dieser Algebra übernimmt der Kommutator die Rolle der Multiplikation. Abgeschlossenheit bedeutet, dass diese Operationen keine Erzeuger außerhalb der Menge hervorbringen.}
\end{classroomqa}

\section{Drehungen um verschiedene Achsen}

Drehungen liefern das wichtigste Beispiel einer nichtabelschen Symmetrie. Eine Drehung um die \(x\)-Achse, gefolgt von einer Drehung um die \(y\)-Achse, ist im Allgemeinen nicht dasselbe wie die Ausführung in umgekehrter Reihenfolge. Die Abweichung hängt selbst mit einer Drehung um die dritte Achse zusammen.

\begin{classroomqa}
\lecturetimestamp{00:54:22}
\audiencequestion{Vertauschen infinitesimale Drehungen miteinander, obwohl endliche Drehungen das nicht tun?}
\lecturerresponse{Sie unterscheiden sich in zweiter Ordnung der beiden kleinen Drehwinkel. Das genügt: Nachdem man diese kleinen Parameter herausgekürzt hat, bleibt ein von null verschiedener Kommutator der Erzeuger selbst.}
\end{classroomqa}

Bevor wir die quantenmechanischen Kommutatoren berechnen, betrachten wir das entsprechende klassische Bild. Nehmen wir eine kreisförmige Planetenbahn mit einem Drehimpuls senkrecht zu ihrer Ebene. Dreht man die gesamte Bahn um die \(y\)-Achse, erhält man eine andere zulässige Bahn mit exakt derselben Energie. Ihr Drehimpulsvektor hat nun jedoch eine andere \(x\)-Komponente.

Rotationsinvarianz erzeugt also Konfigurationen gleicher Energie mit verschiedenen Drehimpulskomponenten. In der Quantenmechanik zeigt sich das als Nichtvertauschbarkeit der Erzeuger von Drehungen um verschiedene Achsen.

\section{Drehimpuls in kartesischen Koordinaten}

Betrachten wir zunächst ein Teilchen, das sich in der \(xy\)-Ebene bewegt, ohne auf einen Kreis beschränkt zu sein. Bei einer kleinen Drehung gegen den Uhrzeigersinn gilt
\begin{equation}
\delta x=-\epsilon y,
\qquad
\delta y=\epsilon x.
\end{equation}
Für den aktiv gedrehten Zustand wird die transformierte Wellenfunktion an den invers gedrehten Koordinaten ausgewertet:
\begin{equation}
\psi_\epsilon(x,y)
=\psi(x+\epsilon y,y-\epsilon x).
\end{equation}
Folglich gilt
\begin{align}
\delta\psi
&=\epsilon\left(
y\frac{\partial}{\partial x}
-x\frac{\partial}{\partial y}
\right)\psi
+\mathcal O(\epsilon^2)\\
&=-\frac{i\epsilon}{\hbar}
\left(xp_y-yp_x\right)\psi
+\mathcal O(\epsilon^2).
\end{align}
Der Erzeuger von Drehungen um die \(z\)-Achse ist daher
\begin{equation}
L_z=xp_y-yp_x.
\end{equation}
Dies ist die \(z\)-Komponente des bekannten klassischen Vektors
\begin{equation}
\mathbf L=\mathbf r\times\mathbf p.
\end{equation}
Sie ist nicht nur der mechanische Drehimpuls; in der Quantenmechanik ist sie der Operator, der Drehungen der Wellenfunktion erzeugt.\footnote{Die Konvention der aktiven Drehung wird so gewählt, dass sie mit der früheren Transformation \(\psi(\theta)\mapsto\psi(\theta-\epsilon)\) und der abschließenden Identifikation \(L_z=xp_y-yp_x\) übereinstimmt, nachdem in der mündlichen Herleitung zwischenzeitlich Unsicherheit über das Vorzeichen aufgetreten war.}

Durch zyklisches Vertauschen von \(x\to y\to z\to x\) erhält man die übrigen Komponenten.\footnote{Die Formeln für \(L_x\) und \(L_y\) erhält man, indem man die in der Vorlesung ausdrücklich vorgegebene zyklische Ersetzung \(x\to y\to z\to x\) auf \(L_z=xp_y-yp_x\) anwendet.}
\begin{align}
L_x&=yp_z-zp_y,\\
L_y&=zp_x-xp_z,\\
L_z&=xp_y-yp_x.
\end{align}

Für ein rotationsinvariantes System vertauschen diese Operatoren mit dem Hamiltonoperator:
\begin{equation}
[H,L_x]=[H,L_y]=[H,L_z]=0.
\end{equation}
Ein Teilchen in einem Zentralkraftfeld ist das Standardbeispiel. Diese Gleichungen sind die Erhaltung des Drehimpulses in quantenmechanischer Form.

\section{Die Lie-Algebra des Drehimpulses}

Die kanonischen Vertauschungsrelationen lauten
\begin{align}
[x_i,x_j]&=0,\\
[p_i,p_j]&=0,\\
[x_i,p_j]&=i\hbar\,\delta_{ij}.
\end{align}
Koordinaten vertauschen mit Impulskomponenten in anderen Richtungen. So vertauscht etwa \(x\) mit \(p_y\), während \(x\) nicht mit \(p_x\) vertauscht.

Berechnen wir nun
\begin{equation}
[L_x,L_y]
=
[yp_z-zp_y,\;zp_x-xp_z].
\end{equation}
Die meisten Terme verschwinden. Die beiden von null verschiedenen Beiträge sind
\begin{align}
[yp_z,zp_x]&=-i\hbar\,yp_x,\\
[zp_y,xp_z]&=i\hbar\,xp_y.
\end{align}
Daher gilt
\begin{align}
[L_x,L_y]
&=i\hbar\left(xp_y-yp_x\right)\\
&=i\hbar L_z.
\end{align}
Durch zyklisches Vertauschen der Koordinatenbezeichnungen erhält man
\begin{align}
[L_x,L_y]&=i\hbar L_z,\\
[L_y,L_z]&=i\hbar L_x,\\
[L_z,L_x]&=i\hbar L_y.
\end{align}

\begin{classroomqa}
\lecturetimestamp{01:14:50}
\audiencequestion{Wenn der Kommutator \(L_z\) ergibt, müssen wir dann noch zeigen, dass \(L_z\) keine Linearkombination von \(L_x\) und \(L_y\) ist?}
\lecturerresponse{Ja. Diese Unabhängigkeit muss nachgewiesen und darf nicht einfach vorausgesetzt werden. Ein direkter algebraischer Beweis wird an dieser Stelle nicht geführt.}
\end{classroomqa}

Die dargestellten Kommutatoren schließen sich auf dem Spann von \(L_x\), \(L_y\) und \(L_z\): Sind alle drei enthalten, erzeugen weitere Kommutatoren keinen Operator außerhalb dieses Spanns. Das ist die Lie-Algebra des Drehimpulses. Zusammen mit
\begin{equation}
[H,L_i]=0,
\qquad i=x,y,z,
\end{equation}
enthalten diese Relationen wesentliche Informationen über das Energiespektrum, ohne dass die Schr\"odinger-Gleichung für die Winkelabhängigkeit direkt gelöst werden muss.

Eine ähnliche algebraische Methode wird beim quantenmechanischen harmonischen Oszillator verwendet, bei dem Erzeugungs- und Vernichtungsoperatoren sein Spektrum erzeugen. Hier verändern die analogen Operatoren die \(z\)-Komponente des Drehimpulses statt der Energie.

Die Wahl der \(z\)-Achse ist lediglich eine Wahl der Darstellung. Nichts macht \(z\) physikalisch vorteilhafter als \(x\) oder \(y\), aber die Wahl einer Komponente liefert eine zweckmäßige Basis. Schreiben wir
\begin{equation}
L_z\ket{m}=m\hbar\ket{m}.
\end{equation}
Zu bestimmen ist, welche Werte von \(m\) auftreten können und wie Zustände mit verschiedenen Werten von \(m\) miteinander zusammenhängen.

\begin{classroomqa}
\lecturetimestamp{01:24:02}
\audiencequestion{Meinen Sie beim Vergleich mit dem harmonischen Oszillator den quantenmechanischen harmonischen Oszillator? Der wurde in der vorherigen Vorlesung nicht behandelt.}
\lecturerresponse{Ja, der Vergleich bezieht sich auf den quantenmechanischen Oszillator. Die vorliegende Argumentation zum Drehimpuls ist in sich abgeschlossen und setzt keine Kenntnis des Oszillators voraus. Die Analogie lässt sich später auch in umgekehrter Richtung betrachten: Beim Oszillator kommt eine ähnliche algebraische Methode zum Einsatz.}
\end{classroomqa}

\section{Drehimpuls anheben und absenken}

Definieren wir
\begin{equation}
L_+=L_x+iL_y,
\qquad
L_-=L_x-iL_y.
\end{equation}
Aus den Drehimpuls-Kommutatoren folgt
\begin{equation}
[L_z,L_+]=\hbar L_+,
\qquad
[L_z,L_-]=-\hbar L_-.
\end{equation}

Wenden wir die erste Relation auf einen Eigenzustand von \(L_z\) an:
\begin{align}
L_zL_+\ket{m}
&=\left(L_+L_z+[L_z,L_+]\right)\ket{m}\\
&=(m+1)\hbar\,L_+\ket{m}.
\end{align}
Sofern \(L_+\ket{m}=0\) nicht gilt, ist dies ein weiterer Eigenzustand von \(L_z\), nun mit dem Eigenwert \((m+1)\hbar\). Entsprechend gilt
\begin{equation}
L_zL_-\ket{m}
=(m-1)\hbar\,L_-\ket{m}.
\end{equation}
Damit erhöht \(L_+\) den Wert von \(m\) um eins und \(L_-\) senkt ihn um eins.

Ausgehend von einem beliebigen Eigenwert erzeugt die Algebra eine Folge mit ganzzahligen Abständen. Diese Folge kann sich unbegrenzt fortsetzen oder enden, weil
\begin{equation}
L_+\ket{m_{\mathrm{top}}}=0
\end{equation}
am oberen Ende und
\begin{equation}
L_-\ket{m_{\mathrm{bottom}}}=0
\end{equation}
am unteren Ende gilt.

\begin{classroomqa}
\lecturetimestamp{01:35:11}
\audiencequestion{War der Drehimpuls durch die Argumentation für den Kreis nicht bereits auf ganzzahlige Werte festgelegt?}
\lecturerresponse{Diese Argumentation weist ganzzahlige Werte für den orbitalen Drehimpuls nach, der durch eindeutige Ortswellenfunktionen dargestellt wird. Die abstrakte Drehimpulsalgebra liefert ein weniger enges Ergebnis: Die Leiter kann ganzzahlige oder halbzahlig Werte enthalten.}
\end{classroomqa}

Endet eine Leiter, verknüpft die Rotationsinvarianz ihre beiden Enden. Eine Drehung um \(180^\circ\) um eine zu \(z\) senkrechte Achse transformiert
\begin{equation}
L_z\longrightarrow-L_z.
\end{equation}
Das Spektrum muss daher symmetrisch um null sein:
\begin{equation}
m_{\mathrm{bottom}}=-m_{\mathrm{top}}.
\end{equation}
Da sich benachbarte Werte um eins unterscheiden, gibt es zwei Möglichkeiten. Die Leiter enthält entweder ganze Zahlen,
\begin{equation}
\ldots,-2,-1,0,1,2,\ldots,
\end{equation}
oder halbe Zahlen,
\begin{equation}
\ldots,-\frac{5}{2},-\frac{3}{2},-\frac{1}{2},
\frac{1}{2},\frac{3}{2},\frac{5}{2},\ldots,
\end{equation}
mit Endpunkten dort, wo der Hebe- beziehungsweise Senkoperator null ergibt.

Die Argumentation mit der eindeutigen Wellenfunktion auf dem Kreis legt für den orbitalen Drehimpuls die ganzzahlige Möglichkeit fest:
\begin{equation}
\mathbf L=\mathbf r\times\mathbf p.
\end{equation}
Auch die halbzahlige Möglichkeit ist physikalisch und gehört zum Spin-Drehimpuls, der gesondert behandelt wird.

Schließlich haben die Leiterzustände dieselbe Energie, wenn der Hamiltonoperator rotationsinvariant ist. Nehmen wir an,
\begin{equation}
H\ket{m}=E\ket{m}.
\end{equation}
Da jede Komponente \(L_i\) mit \(H\) vertauscht, gilt dies auch für \(L_+\) und \(L_-\). Somit gilt
\begin{align}
H L_+\ket{m}
&=L_+H\ket{m}\\
&=E L_+\ket{m}.
\end{align}
Immer wenn \(L_+\ket{m}\neq0\) gilt, hat der angehobene Zustand dieselbe Energie wie der ursprüngliche Zustand. Dieselbe Rechnung ergibt
\begin{equation}
H L_-\ket{m}=E L_-\ket{m}.
\end{equation}
Anheben und Absenken erzeugen somit eine Zustandsfamilie mit verschiedenen Eigenwerten von \(L_z\), aber exakt derselben Energie. Dies ist der Zusammenhang zwischen nichtvertauschenden Symmetrien und Entartung, der sich durch die ganze Vorlesung zieht.

\begin{classroomqa}
\lecturetimestamp{01:44:13}
\audiencequestion{Sind diese Leiterzustände mit der Wahl verschiedener Drehachsen verbunden?}
\lecturerresponse{Sie hängen durch Drehungen zusammen. Eigenvektoren von \(L_x\) oder \(L_y\) sind Linearkombinationen der gewählten \(L_z\)-Eigenvektoren. Eine Drehung des Zustands verändert seine Drehimpulskomponenten, lässt seine Energie aber unverändert.}
\end{classroomqa}

Die Wahl von \(L_+\) und \(L_-\) ist zweckmäßig, weil diese Kombinationen die Leiterstruktur unmittelbar sichtbar machen. Dieser Kunstgriff tritt in der Quantenmechanik immer wieder auf: Schließt sich eine Kommutatoralgebra, ordnen Hebe- und Senkoperatoren häufig ihr Eigenwertspektrum.

\begin{classroomqa}
\lecturetimestamp{01:46:33}
\audiencequestion{Was geschieht, wenn \(L_+\) auf den Zustand mit dem maximalen Wert wirkt, und wie ist der oberste Zustand definiert?}
\lecturerresponse{Das Ergebnis ist null. Das ist die Definition des obersten Zustands einer endenden Leiter. Entsprechend ergibt \(L_-\) am unteren Ende null. Eine abstrakte Leiter muss nicht enden. Bei den später betrachteten atomaren Drehimpulszuständen trifft man typischerweise auf endende Leitern.}
\end{classroomqa}

Als Nächstes entwickeln wir diese Drehimpulsstruktur weiter und wenden sie auf Atome einschließlich des Wasserstoffatoms an. Die hier gefundenen Rotationsmultipletts erklären einen Teil des Entartungsmusters; weitere Strukturen müssen noch entwickelt werden.
